\begin{chaptersummary}
  \item A subgraph is a package reference, one call on the builder, and two
        attributes per entity. Everything else the contract asks for is
        generated: \code{@link}, the directive definitions, \code{\_service},
        \code{\_entities}, the \code{\_Entity} union and \code{\_Any} are all
        written into the exported schema by the package, and the union grows
        and shrinks with the number of types carrying a key.
  \item \code{\_service} hands back the schema as a string, and it is the same
        document \code{schema export} writes rather than a federation-only view
        of it. The specification says the field is intended solely for the
        router; it is nonetheless on \code{Query}, it introspects, and anything
        that can reach the service can ask for it.
  \item A reference resolver is one substep of the specification's algorithm and
        nothing else: the package owns the loop, the ordering and the dispatch
        on \code{\_\_typename}. What it hands you is the representation as it
        arrived, which for this book means the global node id as an undecoded
        string. \code{[ID<T>]} does not decode it and does not complain; it
        yields zero and the entity comes back null.
  \item The entity route is not \code{node(id:)} with a different spelling. It
        takes a different argument, it cannot be given a
        \code{QueryContext<T>}, and without one the projection is gone: six
        columns come back where the node route asked for one. Whether a batch
        of representations costs one statement or one each is decided by what
        is behind your resolver, which is the whole of
        chapter~\ref{ch:router-n-plus-1} arriving early.
  \item A key is a value another service produced, and the only thing that
        checks it is the code you write. A node id naming one type resolves
        perfectly well against another, because both decode to an integer, and
        the answer is well typed, correctly shaped, HTTP 200 and wrong. Compare
        the type name inside the key before spending the integer.
\end{chaptersummary}
